% 344_Furey_GALG_FFGFT_En_ch.tex
% Johann Pascher, 3 September 2026

\providecommand{\BM}{\textbf{[B]}}
\providecommand{\KM}{\textbf{[K]}}
\providecommand{\SM}{\textbf{[S]}}

\chapter*{Furey Fermions in GALG and FFGFT: N-Operator, Tripotents, and Neutrino-Vacuum Identity}
\addcontentsline{toc}{chapter}{Furey Fermions in GALG and FFGFT}

\section*{Abstract}

Doug Matzke's IPI mail of 3~September 2026 presents concrete GALG results on Furey fermions
in Cl$(6)$: complete Up/Down fermion sets Su/Sd, idempotents $N_1,N_2,N_3$, tripotent $N$,
jE7 projector, and a grade-1+grade-5 construction. The strongest single result:
$\mathrm{Su}[0] = \mathrm{Vss}[0]$ exactly --- the neutrino is the lowest vacuum state \BM.
Theorem~F shows: the tripotent property for even vacua and the Frobenius fixed-point property
(Doc.~339) are the same decomposition from two independent directions \BM.
The electron exception at the jE7 projector is the known Furey problem \SM.

\section{Theorem A --- $N_k$ are idempotents in characteristic 3}

$N_k = -1 + (-i\,B_k)$, $B_k^2 = -1$ (Witt pair):
\[
N_k^2 = 1 + 2i\,B_k + (-1)(-1) = 2 + 2i\,B_k \equiv -1 + (-i\,B_k) \pmod{3} = N_k. \quad \BM
\]
The sum loses its scalar part since $-3\equiv0$: $N_1+N_2+N_3 \equiv (-i)\sum B_k = N$.

\section{Theorem B --- $N$ is a tripotent}

Cross-terms $N_kN_j$ ($k\neq j$) are pure grade-4 elements; hence $N^3=N$ \BM.
Doug's output: \texttt{N**3 = N}, \texttt{N**4 = N**2}.

\section{Theorem C --- $\mathrm{Su}[0] = \mathrm{Vss}[0]$: neutrino is the vacuum}

\texttt{Su[0] == Vss[0]: True} and \texttt{(1+N)*(1+jE7) == Vss[0]: True} \BM.

\textbf{FFGFT correspondence \BM.} In FFGFT, $\nu_1$ lies in orbit $O_1=\{1,3,9\}$, the orbit
of the generator of GF$(27)^*$ --- the algebraically simplest object. Both frameworks
independently place the neutrino at the algebraically simplest position: convergence of
two independent frameworks.

\section{Theorems D--F}

\textbf{D --- sigma \BM.} $e_1\to e_2\to e_4\to e_1$, $e_3\to e_6\to e_5\to e_3$;
\texttt{\{1:2,2:4,4:1,3:6,6:5,5:3\}}, two 3-cycles, $\sigma^3=\mathrm{id}$.
Doc.~341 contains this formula; Doug's \texttt{all\_g1\_g5good} is the $\sigma$-invariant sum.

\textbf{E --- jE7 $\equiv$ $N_k$ as projectors \BM; electron exception \SM.}
$F\cdot jE7 = +F$ for all Su (8/8); $F\cdot jE7 = -F$ for 7/8 of Sd (Sd[7]=electron fails).
Same for $N_k$. Known Furey problem. FFGFT: electron in Dirac layer $\{f_3,f_4\}$,
neutrino in Majorana layer $\{f_1,f_2\}$ (Doc.~343, Theorem~D$'''$). Common root open \SM.

\textbf{F --- Tripotent for even vacua \BM.}
$(V_k+P)^2=-V_k$, $(V_k+P)^3=V_k$ for $k\in\{0,2,4\}$;
no collapse for $k\in\{1,3,5\}$ (length 32). Identical to the Frobenius
fixed-point/orbit split of Doc.~339 --- two independent derivations of the same decomposition.

\section{Summary}

\begin{table}[H]\centering
\caption{Results Doc.~344}
\begin{tabular}{lll}\toprule
Theorem & Content & Status \\\midrule
A & $N_k^2=N_k$, characteristic-3 proof & \BM \\
B & $N^3=N$, tripotent & \BM \\
C & Su[0]=Vss[0]: neutrino=vacuum; FFGFT correspondence & \BM \\
D & sigma$=\{1\to2,2\to4,4\to1,3\to6,6\to5,5\to3\}$ & \BM \\
E & jE7$\equiv N_k$; electron exception (Furey problem) & \BM/\SM \\
F & Tripotent even vacua = Frobenius fixed points (Doc.~339) & \BM \\\bottomrule
\end{tabular}
\end{table}

\textbf{Open question \SM:} Common algebraic root of the electron inconsistency in GALG/Furey
and the Majorana/Dirac layer separation in FFGFT (Doc.~343, Theorem~D$'''$).

\section*{Verification script}
\texttt{pruef\_344\_furey\_galg.py} --- 6 theorems, all assertions passed.

\begin{thebibliography}{9}
\bibitem{Furey2016} C.~Furey, \textit{Standard Model Physics from an Algebra?} (2016).
\bibitem{Matzke2026} D.~Matzke, IPI mail 3~September 2026.
\bibitem{P339} J.~Pascher, \textit{Doc.~339: Frobenius separation}, FFGFT corpus (Aug. 2026).
\bibitem{P341} J.~Pascher, \textit{Doc.~341: GF(27) in GALG and FFGFT}, FFGFT corpus (Sept. 2026).
\bibitem{P343} J.~Pascher, \textit{Doc.~343: Zeta function of T$^4$/Z$_3$}, FFGFT corpus (Sept. 2026).
\end{thebibliography}
